\documentclass[11pt]{article}

% ---------- packages ----------
\usepackage[utf8]{inputenc}
\usepackage[T1]{fontenc}
\usepackage{lmodern}
\usepackage[margin=1in]{geometry}
\usepackage{amsmath, amssymb}
\usepackage{amsthm}
\usepackage{booktabs}
\usepackage{enumitem}
\usepackage{microtype}
\usepackage{indentfirst}
\usepackage{float}
\usepackage[hidelinks]{hyperref}

% ---------- theorem-like environments ----------
\newtheorem{theorem}{Theorem}
\newtheorem{proposition}{Proposition}
\newtheorem{lemma}{Lemma}
\theoremstyle{remark}
\newtheorem{remark}{Remark}

% ---------- shorthands (all explained in the Notation section) ----------
\newcommand{\Prob}{\operatorname{P}}
\newcommand{\Bin}{\operatorname{Binomial}}
\newcommand{\Mult}{\operatorname{Multinomial}}
\newcommand{\Pois}{\operatorname{Poisson}}
\newcommand{\E}{\operatorname{E}}

\title{\bfseries Perfecting a Final Chocobo:\\
Exact Probabilities and Optimal Feeding Policies\\
for \emph{Final Fantasy XIV} Chocobo Racing}
\author{Kevin Rutledge}
\date{June 2026}

\begin{document}
\maketitle

\begin{abstract}
In \emph{Final Fantasy XIV}, a fully bred ``final'' race chocobo caps all five
attributes at $500$ points. Its lifetime growth gives $245$ random increments and
$50$ player-directed feedings, and that total cannot fill all five attributes. Some
attribute must be sacrificed, conventionally \emph{Acceleration}. I compute the
probability of perfecting the other four under two regimes of play. In the
\textbf{fixed-end} regime the player feeds only at rank $50$. I prove that the
popular $6.17\%$ community figure is the answer to a frictionless idealization,
compute the \emph{exact} probability once feeds are treated as discrete
$5/10/15$-point chunks ($2.30\%$ for a pre-committed dump and $11.32\%$ when the
dump is chosen after the growth is seen), and show that perfecting \emph{some} three
attributes is certain. In the \textbf{adaptive} regime the player feeds
\emph{during} the climb and exploits the rule that a maxed attribute leaves the
random pool. I cast this as a Markov decision process and solve it by exact
backward induction. The best \emph{online} probability is what a player can achieve
deciding rank by rank, without foreknowledge of the rolls. For the strict
Grade-3-only lineup $500/500/500/500/250$ it falls between $1.06\%$ and $1.17\%$,
about $20\%$ above the best fixed-deadline rule. I bracket the \emph{offline}
ceiling, what a player who could see every roll in advance could reach, exactly
within $[0.10\%,\,6.17\%]$, and pin it at $4.48\%$ by high-precision simulation. The
closed-form and backward-induction results are exact. Three adaptive-regime
quantities are Monte-Carlo estimates, flagged where they appear. The paper is
self-contained, and every probabilistic tool is defined where it is first used.
\end{abstract}

\tableofcontents
\bigskip

% =====================================================================
\section{Introduction}\label{sec:intro}
% =====================================================================

\emph{Chocobo racing} is a long-running minigame in \emph{Final Fantasy XIV}. A
player raises a bird through fifty ranks, and its performance on the track is
governed by five numerical attributes, namely \emph{Maximum Speed},
\emph{Acceleration}, \emph{Endurance}, \emph{Stamina}, and \emph{Cunning}. Through
breeding each attribute acquires a \emph{cap}, and a fully bred ``final'' chocobo,
the endpoint of the breeding game, caps all five at the maximum of $500$ points.
This paper asks one quantitative question about such a bird. \emph{How likely is it
to finish with four of its attributes exactly at the cap?}

The question is interesting precisely because the answer is far below one. An
attribute grows in two ways. It receives a little \emph{random} growth at every
rank-up, and it receives \emph{feeding}, in which the player spends a limited supply
of food to add points to a chosen attribute. As the next sections make precise, the
lifetime totals of $245$ random increments and $50$ feedings are together not enough
to lift all five attributes to $500$. The player must let at least one attribute
fall short, and the natural sacrifice is \emph{Acceleration}, the attribute that
affects race results least and so serves as the conventional ``dump.'' The
aspirational lineup is therefore $500/500/500/500/250$.

How likely the four-cap finish is depends sharply on \emph{how} the player feeds,
and I study two regimes. In the \textbf{fixed-end} regime the player banks every
feeding and spends it all at rank $50$. Here I explain the $6.17\%$ figure quoted
in community guides, show that it answers only a frictionless idealization, and
compute the true probability once feeds are recognized as discrete $5$-, $10$-, and
$15$-point portions. I also settle the easier goal of perfecting merely three
attributes. In the \textbf{adaptive} regime the player feeds \emph{during} the
climb and exploits a subtle rule to steer the outcome. A maxed attribute stops
receiving random growth. There the question becomes one of optimal control, and the
quantities of interest are the best achievable success rate and the ceiling that a
player able to foresee the rolls could reach.

The paper has a shared setup followed by two parts. The setup fixes the game model,
the elementary probabilistic tools, and the model of a single climb. \textbf{Part~I}
treats fixed-end feeding. It covers the idealized $6.17\%$ bound, the exact chunked
probabilities ($2.30\%$ for a pre-committed dump and $11.32\%$ for a dump chosen
after the fact), and the certainty of a three-attribute finish. \textbf{Part~II}
treats adaptive feeding. It builds a Markov decision process whose exact solution
brackets the online optimum of the strict Grade-3-only lineup, gives an explicit
policy that realizes it, and proves exact two-sided bounds on the offline ceiling. A
closing section collects every result in a single table and translates the optimal
policy into a rule a player can follow. Results obtained in closed form or by exact
backward induction are labeled \emph{exact}. The few obtained by simulation are
labeled \emph{Monte-Carlo} wherever they appear.

% =====================================================================
\section{The game model}\label{sec:model}
% =====================================================================

The in-game rules translate into arithmetic. The rules below are as documented
in the community wikis~\cite{cgw,wp-rankup,wp-final,wp-train,icy}. They are a fixed model with no
modeling approximations beyond those noted.

\paragraph{Attributes and caps.}
A race chocobo has five attributes. Breeding raises their \emph{caps}, and for the
final chocobo studied here every cap equals the maximum, $500$ points. The five caps
are fixed at $500$ throughout.

\paragraph{The percent-to-point conversion.}
The game expresses all growth as a percentage of an attribute's cap. At a cap of
$500$,
\[
   1\%\text{ of cap} \;=\; 5\text{ points.}
\]
Because $500$ is a multiple of $100$, every percentage quoted below is a whole
number of points, and there is no rounding to track. (At a different cap, $1\%$ need
not be an integer, and rounding would matter. The integer-point arithmetic of this
paper is special to the cap-$500$ final chocobo.)

\paragraph{Starting values.}
A freshly hatched rank-$1$ chocobo begins with every attribute at $11\%$ of its cap,
that is, at $55$ points.

\paragraph{Random growth from ranking up.}
Reaching rank $50$ from rank $1$ takes $49$ rank-ups. \emph{Each} rank-up applies
five separate $+1\%$ ($+5$-point) \emph{increments}. Each increment is placed on one
attribute chosen uniformly at random and independently of the others. Over a full
climb the number of increments is therefore
\[
   N \;=\; 49\times 5 \;=\; 245.
\]

\paragraph{Feeding.}
The chocobo holds one feeding \emph{slot} at rank $1$ and gains one more at each
rank-up, for $50$ slots over its lifetime. Unused slots are stockpiled, so the
player chooses freely \emph{when} to spend them. Each slot holds a single feeding of
grade $1$, $2$, or $3$, which adds the points in Table~\ref{tab:feed} to one chosen
attribute. A feeding may not raise an attribute above its $500$ cap, and any
overshoot is wasted, so to \emph{max} an attribute the player must land on $500$
exactly.

\begin{table}[h]
\centering
\begin{tabular}{@{}cccl@{}}
\toprule
Feed grade & Growth (\% of cap) & Points added (cap $500$) & Slots used \\
\midrule
Grade 1 & $1\%$ & $+5$  & $1$ \\
Grade 2 & $2\%$ & $+10$ & $1$ \\
Grade 3 & $3\%$ & $+15$ & $1$ \\
\bottomrule
\end{tabular}
\caption{The three feed grades. Every feeding occupies exactly one of the $50$
lifetime slots, regardless of grade, so a higher grade delivers more points per
slot. Grade-3 feeding is the most slot-efficient and is central to Part~II.}
\label{tab:feed}
\end{table}

% =====================================================================
\section{Notation and elementary tools}\label{sec:notation}
% =====================================================================

The mathematics is kept elementary, and the following objects are all that the
arguments require. Two heavier tools are introduced in full where they are first
needed, namely \emph{generating functions} (used in Part~I) and \emph{Markov
decision processes} solved by \emph{backward induction} (used in Part~II).

\begin{description}[leftmargin=2.2em, style=nextline]
\item[Factorial $n!$] the product $1\cdot 2\cdots n$, with $0!=1$. It counts the
number of orderings of $n$ distinct items.

\item[Binomial coefficient $\binom{n}{k}$] the number of ways to choose $k$ items
from $n$, equal to $\dfrac{n!}{k!\,(n-k)!}$.

\item[Ceiling $\lceil x\rceil$] the smallest whole number that is not below $x$.
For example $\lceil 3.0\rceil=3$ and $\lceil 3.2\rceil=4$. It counts
``how many feedings are needed,'' since a partial feeding still consumes a whole
slot.

\item[Modular residue $n \bmod 3$] the remainder of $n$ on division by $3$, one of
$0,1,2$. It tracks whether an attribute's gap to $500$ is divisible by $15$ and so
fillable by Grade-3 feeds alone, the pivotal condition of Part~II.

\item[Probability $\Prob(E)$] a number in $[0,1]$ measuring how often event $E$
occurs. $\E[X]$ denotes the average value of a quantity $X$.

\item[Binomial distribution $\Bin(m,p)$] the random count of ``successes'' in $m$
independent yes/no trials, each succeeding with probability $p$. The chance of
exactly $k$ successes is
\[
   \Prob(\text{count}=k)=\binom{m}{k}\,p^{k}(1-p)^{\,m-k},
\]
and the average count is $m\,p$.

\item[Multinomial distribution $\Mult(m;\,p_1,\dots,p_r)$] the generalization to
$r$ categories. Dropping $m$ balls independently into $r$ boxes, box $j$ chosen
with probability $p_j$, gives box counts $(n_1,\dots,n_r)$ with
\[
   \Prob(n_1,\dots,n_r)=\frac{m!}{n_1!\,n_2!\cdots n_r!}\;p_1^{n_1}\cdots p_r^{n_r}.
\]
Each individual box count is itself $\Bin(m,p_j)$.

\item[Poisson distribution $\Pois(\lambda)$] the count of rare independent events
with average $\lambda$. The chance of exactly $n$ is $e^{-\lambda}\lambda^{n}/n!$.
One convenient fact about it, called \emph{Poissonization}, is used in Part~I and
explained there.
\end{description}

% =====================================================================
\section{The probabilistic model of a climb}\label{sec:prob}
% =====================================================================

Label the five attributes $0,1,2,3,4$ and let attribute $0$ be Acceleration, the
dump attribute. Write $n_i$ for the number of random increments that land on
attribute $i$ during the climb.

\begin{lemma}[Increment counts]\label{lem:mult}
If no attribute is maxed before rank $50$, then the counts $(n_0,\dots,n_4)$ follow
$\Mult\!\left(245;\,\tfrac15,\dots,\tfrac15\right)$. In particular each $n_i$ is
$\Bin(245,\tfrac15)$ with average $49$.
\end{lemma}

\begin{proof}
Each of the $N=245$ increments is placed, independently, on a uniformly chosen
attribute, exactly the ``balls into boxes'' description of the multinomial law,
with $245$ balls and five equally likely boxes. The lone caveat is that the game
removes an attribute from the pool once it is maxed. From random growth alone, an
attribute maxes only on reaching $55+5n_i=500$, that is $n_i=89$ increments. The
chance that any attribute does so is at most $1.2\times10^{-8}$ (the bound is
computed in Part~I), so under the stated hypothesis the five boxes stay equally
likely throughout.
\end{proof}

\begin{remark}
The hypothesis ``no attribute maxed early'' holds automatically in the fixed-end
regime of Part~I, where nothing is fed until rank $50$. Part~II deliberately maxes
attributes \emph{during} the climb (``locking''), which triggers the exclusion the
Lemma sets aside. Its analysis accounts for that effect directly.
\end{remark}

\paragraph{Deficits.}
Before feeding, attribute $i$ sits at $55+5n_i$ points, so its gap to the cap is
\begin{equation}\label{eq:deficit}
   D_i \;=\; 500-(55+5n_i) \;=\; 5\,(89-n_i)\quad\text{points.}
\end{equation}

\paragraph{Feeding cost.}
Since $D_i$ is a multiple of $5$, the fewest slots that close the gap \emph{exactly}
use as many $15$-point (Grade-3) feeds as possible and one smaller feed for any
remainder. The minimum number of slots to max attribute $i$ is therefore
\begin{equation}\label{eq:cost}
   c(n_i) \;=\; \Big\lceil \tfrac{D_i}{15}\Big\rceil
            \;=\; \Big\lceil \tfrac{89-n_i}{3}\Big\rceil
   \qquad(\text{and }0\text{ once }n_i\ge 89).
\end{equation}
The ceiling is the ``chunkiness'' that the idealized $6.17\%$ estimate of Part~I
overlooks, since a leftover of $5$ or $10$ points still costs a whole slot.

\paragraph{Success condition.}
With all $50$ slots in hand at rank $50$, maxing a chosen set $T$ of attributes to
exactly $500$ succeeds precisely when their costs fit the budget,
\begin{equation}\label{eq:success}
   \sum_{i\in T} c(n_i) \;\le\; 50.
\end{equation}
(Spending every slot at rank $50$ makes slot \emph{timing} automatic in the
fixed-end regime. Part~II reinstates timing as a genuine constraint.)

\paragraph{The goals analyzed.}
Table~\ref{tab:goals} names the four scenarios studied in this paper and points to
where each is treated. They differ along three axes, namely how many attributes are
maxed, whether the sacrificed attribute is fixed in advance or chosen after the
growth is seen, and whether any feed grade may be used or only the slot-efficient
Grade-3.

\begin{table}[h]
\centering
\begin{tabular}{@{}lcccl@{}}
\toprule
Goal & Maxed & Dump & Feed grades & Regime \\
\midrule
Four, fixed dump   & $4$ & Acceleration (pre-set) & any        & fixed-end (Part~I) \\
Four, flexible dump& $4$ & worst attribute        & any        & fixed-end (Part~I) \\
Three              & $3$ & two worst              & any        & fixed-end (Part~I) \\
Perfect            & $4$ & Acceleration (pre-set) & Grade-3 only & adaptive (Part~II) \\
\bottomrule
\end{tabular}
\caption{The four goals. The first three (fixed-end) ask only that four or three
attributes reach $500$ using feeds of any grade. The fourth (adaptive) is stricter.
It insists on the exact lineup $500/500/500/500/250$ reached with Grade-3 feeds
alone, with no partial feed ever wasted, and it lets the player feed during the climb.}
\label{tab:goals}
\end{table}

% =====================================================================
\part{Fixed-end feeding}
% =====================================================================

In this part the player banks every feeding and spends it all at rank $50$. I
first show this schedule is optimal for the goals at hand
(Section~\ref{sec:feedlast}), then read off the resulting probabilities. These are
the idealized $6.17\%$ that explains the community figure (Section~\ref{sec:naive}),
the exact chunked values $2.30\%$ and $11.32\%$ (Section~\ref{sec:exact}), and the
certainty of a three-attribute finish (Section~\ref{sec:three}).

% =====================================================================
\section{Feeding last is optimal}\label{sec:feedlast}
% =====================================================================

\begin{proposition}[Optimality of feeding last]\label{prop:feedlast}
For each fixed-end goal of Table~\ref{tab:goals}, which maxes a chosen set of
attributes using feeds of any grade, storing every feeding until rank $50$ maximizes
the probability of success.
\end{proposition}

\begin{proof}
The increments are dealt by chance, and the player controls only \emph{when} to
feed. Feeding can change the increment distribution in just one way. Maxing an
attribute early, by Lemma~\ref{lem:mult}, drops it from the pool, so that every
later increment falls on the attributes that remain. For any such goal, success
depends through \eqref{eq:cost}--\eqref{eq:success} on the targets' final deficits,
and a deficit shrinks only when its attribute absorbs increments. Maxing a target
early freezes it and reroutes its remaining increments onto the other attributes,
which enlarges their deficits or pads a dump, while maxing a dump early is never
useful. Hence no early feeding lowers the required budget, and some raise it. Feeding
last performs no early maxing, so it minimizes the budget needed. It also leaves the
increments exactly multinomial (Lemma~\ref{lem:mult}).
\end{proof}

\begin{remark}
The proof relies on the fact that, with mixed feed grades, \emph{any} deficit can be
closed exactly at rank $50$. Part~II abandons that freedom. Using Grade-3 feeds
only, a target is maxable solely on a divisible-by-$15$ window, and maxing one early
can \emph{create} a window that rank-$50$ play would miss. That benefit is absent
here, and it is exactly why adaptive feeding overtakes feed-last for the
Grade-3-only goal. The optimality above is therefore specific to the any-grade goals
of Part~I and does not hold in general.
\end{remark}

% =====================================================================
\section{The idealized bound: where \texorpdfstring{$6.17\%$}{6.17\%} comes from}\label{sec:naive}
% =====================================================================

Suppose for a moment that feedings were infinitely divisible, able to deliver any
number of points rather than only the $5$, $10$, or $15$ of a real feed. Closing a
gap of $D_i$ points would
then cost $D_i/15$ slots with no ceiling. Take the ``four, fixed dump'' goal,
$T=\{1,2,3,4\}$, and write $a=n_0$ for Acceleration's count, so the four targets
share the remaining $245-a$ increments. Using the deficit \eqref{eq:deficit} and the
identity $\sum_{i=1}^{4} n_i = 245-a$, the idealized cost of the four targets is
\[
   \sum_{i=1}^{4}\frac{D_i}{15}
   = \frac{5\big(4\cdot 89-(245-a)\big)}{15}
   = \frac{555+5a}{15}
   = 37+\frac{a}{3}.
\]
Requiring this to be at most $50$ gives the clean condition
\begin{equation}\label{eq:naivecond}
   37+\frac{a}{3}\le 50 \iff a\le 39.
\end{equation}
So the idealized problem succeeds exactly when Acceleration draws at most $39$ of the
$245$ increments. By Lemma~\ref{lem:mult}, $a\sim\Bin(245,\tfrac15)$, so this
probability is the exact finite sum
\begin{equation}\label{eq:binomcdf}
   \Prob(a\le 39)=\sum_{k=0}^{39}\binom{245}{k}\Big(\tfrac15\Big)^{k}\Big(\tfrac45\Big)^{245-k}.
\end{equation}

\begin{proposition}\label{prop:naive}
The sum \eqref{eq:binomcdf} equals $6.166840\%$ (to six decimal places).
\end{proposition}

Expression \eqref{eq:binomcdf} is a closed form, a finite sum of whole numbers over
$5^{245}$, evaluated in exact integer arithmetic. Its value, $6.166840\%$,
is the figure community guides round to ``about $6.17\%$.'' It is the right answer to
the \emph{idealized} problem. Note that the average of $a$ is $49$, comfortably above
the threshold $39$, so success is a genuine ``lucky low roll,'' an event well into
the lower tail.

\begin{remark}
The same sum supplies the bound promised in Lemma~\ref{lem:mult}. One particular
attribute reaches $89$ increments with probability
$\sum_{k=89}^{245}\binom{245}{k}(\tfrac15)^k(\tfrac45)^{245-k}\approx 2.3\times
10^{-9}$, and at most five attributes can, giving the $1.2\times 10^{-8}$ union
bound quoted there.
\end{remark}

% =====================================================================
\section{The exact analysis with chunked feeding}\label{sec:exact}
% =====================================================================

The idealized bound overcounts successes because real feeds cannot deliver a
fractional slot's worth of points. The honest success condition keeps the ceiling,
that is \eqref{eq:success} with $c$ from \eqref{eq:cost}. I now evaluate it exactly.

\subsection{A counting (generating-function) approach}
The exact value sums the multinomial probability of \emph{every} increment outcome
$(n_0,\dots,n_4)$ that satisfies the success condition. These are far too many to
list, so a standard bookkeeping device handles them, the \emph{generating function}, a
polynomial whose exponents store the tracked quantities. With two placeholder
symbols $x$ and $y$,
the exponent of $x$ records how many increments an attribute used and the exponent of
$y$ records how many slots it will cost. For one attribute,
\begin{equation}\label{eq:gf}
   g(x,y)\;=\;\sum_{n\ge 0} w(n)\,x^{\,n}\,y^{\,c(n)},
\end{equation}
where $w(n)$ is the probabilistic weight of receiving $n$ increments and $c(n)$ is
the feeding cost \eqref{eq:cost}. Multiplying generating functions \emph{adds} their
exponents, so multiplying four copies of $g$ enumerates every way the four targets
can split their increments while summing the costs. Reading off the coefficient of
$x^{245}$ pins the total increment count to its true value, and keeping only terms of
$y$-exponent at most $50$ enforces the budget \eqref{eq:success}. The bookkeeping is
exact. It merely organizes a vast sum so a computer can evaluate it term by term.

\subsection{Making the weights independent: Poissonization}
One obstacle remains. In the multinomial law the five counts must add to $245$, so the
attributes are dependent and their generating functions do not cleanly multiply. A
classical trick, \emph{Poissonization}, removes the link, resting on one fact.

\begin{quote}
If five independent $\Pois(49)$ counts are observed and only the outcomes that
happen to sum to $245$ are kept, the resulting counts follow exactly
$\Mult(245;\tfrac15,\dots,\tfrac15)$.
\end{quote}

Independent Poisson counts ``do not know'' about one another, so they multiply
freely, and conditioning on their sum afterwards rebuilds the desired multinomial.
Concretely, for
every event $S$ on the counts,
\begin{equation}\label{eq:poiss}
   \Prob(S)\;=\;\frac{\Prob\big(S\ \text{and}\ \textstyle\sum_i N_i=245\big)}
                     {\Prob\big(\textstyle\sum_i N_i=245\big)},
   \qquad N_i\stackrel{\text{ind.}}{\sim}\Pois(49).
\end{equation}
The denominator is a single value. A sum of independent $\Pois(49)$ counts is
$\Pois(245)$, so $\Prob(\sum_i N_i=245)=e^{-245}245^{245}/245!$. In the numerator the
attributes are independent, so $w(n)=e^{-49}49^{n}/n!$ is used in \eqref{eq:gf}, and
the generating-function product applies directly.

\subsection{Evaluation and results}
For the four-fixed-dump goal the computation forms $g(x,y)^4$ (the four targets), multiplies by
$h(x)=\sum_n w(n)x^n$ for the unconstrained Acceleration attribute, extracts the
coefficient of $x^{245}$, sums the terms of $y$-exponent at most $50$, and divides by the
denominator of \eqref{eq:poiss}. This finite computation was carried out exactly,
organized as a table of running coefficients. Two internal checks confirm it. The
total probability over \emph{all} outcomes returns $1.0000000000$, and replacing the
ceiling \eqref{eq:cost} by the idealized cost reproduces Proposition~\ref{prop:naive}'s
$6.166840\%$ to the digit.

\begin{theorem}[Four, fixed dump]\label{thm:main}
Under feed-last play, with Acceleration pre-committed as the dump, the probability of
maxing the other four attributes to exactly $500$ is
\[
   \boxed{\,2.301513\%\,}.
\]
\end{theorem}

The \emph{flexible-dump} goal maxes any four attributes, discarding whichever proves
most expensive. It uses the same machinery with the bookkeeping extended to also track
the single largest per-attribute cost, so the worst attribute can be dropped after the
growth is seen. Its probability is $11.324144\%$. The gap between $2.30\%$ and
$11.32\%$ is the price of committing to Acceleration in advance rather than sacrificing
whatever rolls worst.

% =====================================================================
\section{Maximizing three attributes is certain}\label{sec:three}
% =====================================================================

\begin{theorem}[Three attributes]\label{thm:three}
Under feed-last play, the three highest-count attributes can always be maxed, and the
three-attribute goal succeeds with probability exactly $1$.
\end{theorem}

\begin{proof}
Sort the counts $n_{(1)}\ge n_{(2)}\ge n_{(3)}\ge n_{(4)}\ge n_{(5)}$. The two
smallest are each at most the average $49$, so $n_{(4)}+n_{(5)}\le \tfrac{2}{5}\cdot
245 = 98$, and therefore
\[
   n_{(1)}+n_{(2)}+n_{(3)}\;\ge\;245-98\;=\;147.
\]
By the deficit formula \eqref{eq:deficit}, the combined deficit of the three largest
attributes is at most
\[
   \sum_{\text{top }3} D_i
   = 5\big(3\cdot 89-(n_{(1)}+n_{(2)}+n_{(3)})\big)
   \le 5\,(267-147) = 600\ \text{points.}
\]
Each ceiling obeys $\lceil D_i/15\rceil \le D_i/15 + \tfrac{14}{15}$, so the total
feeding cost is bounded by
\[
   \sum_{\text{top }3}\Big\lceil \tfrac{D_i}{15}\Big\rceil
   \;\le\; \frac{600}{15}+3\cdot\frac{14}{15}
   \;=\;40+2.8\;=\;42.8,
\]
hence at most $42$ slots, an integer. Since $42\le 50$, the budget always suffices,
whatever the random outcome.
\end{proof}

This is why three-attribute builds are treated as routine in practice. They are
guaranteed, well beyond merely likely. The harder goal of maxing three \emph{specific}
pre-chosen attributes, where the two dump attributes may hoard increments, is only
near-certain rather than certain, and is not analyzed here.

% =====================================================================
\part{Adaptive feeding}
% =====================================================================

In this part the player feeds \emph{during} the climb rather than only at rank $50$.
That freedom unlocks a stricter target and turns the analysis into one of optimal
control. This part describes the opportunity (Section~\ref{sec:adaptive}), casts the climb as a
Markov decision process and solves it for the best \emph{online} policy
(Section~\ref{sec:mdp}), states the explicit policy a player can follow
(Section~\ref{sec:policy}), bounds the \emph{offline} ceiling exactly
(Section~\ref{sec:ceiling}), and examines the stricter exact lineup
$500/500/500/500/250$ (Section~\ref{sec:exact250}).

% =====================================================================
\section{The adaptive opportunity}\label{sec:adaptive}
% =====================================================================

\paragraph{A stricter goal.}
Call a finish \emph{perfect} if the four targets reach exactly $500$ using
\emph{only} Grade-3 feeds, with no $5$- or $10$-point top-off ever spent, landing the
lineup $500/500/500/500/250$. By the cost formula \eqref{eq:cost}, a target is
maxable with Grade-3 feeds alone only when its count satisfies $n_i\equiv 2\pmod 3$
at the moment it is fed (then its deficit $D_i=5(89-n_i)$ is a multiple of $15$). This
is strictly harder than the any-grade goal of Part~I, and the surcharge for ``no
wasted feed'' is exactly this divisibility.

\paragraph{A smarter move.}
As an attribute absorbs increments its count cycles through residues
$0\to1\to2\to0\to\cdots$, so a \emph{window} ($n_i\equiv 2$) reopens every three
increments. Rather than gamble that a target lands on a window at rank $50$, the
player can \emph{lock} it mid-climb and feed it to $500$ at a window. Locking carries
a cost. By Lemma~\ref{lem:mult} a maxed attribute leaves the random pool, so every later
rank-up's five increments fall on the attributes that remain, Acceleration included.
Locking early therefore trades ``catch a good window now'' against ``inflate the dump
later,'' and deciding when to lock each target is a sequential decision problem.

\paragraph{Online versus offline.}
Two yardsticks bracket any policy. An \emph{offline} player who could see every
future roll would lock each target at the last window that keeps Acceleration in
budget, and this ceiling is $\approx 4.48\%$, bounded exactly in
Section~\ref{sec:ceiling}. A real \emph{online} player decides from the current state
alone and does worse. The distance between them is the \emph{prophet gap}, the value
of information, which no online rule can recover~\cite{prophet}.

% =====================================================================
\section{The decision process and a budget identity}\label{sec:mdp}
% =====================================================================

The climb is modeled as a finite-horizon \emph{Markov decision process}. At each rank-up
the five increments land (Lemma~\ref{lem:mult}), then the player may lock any target
sitting on a window. The best \emph{online} policy is found by \emph{backward
induction}. Working from the final rank-up backward, the best action in each state
solves $V_t(s)=\max_a\{\,r_t(s,a)+\E[V_{t+1}]\,\}$~\cite{mdp}.

The state looks intractable. Four target counts (each $0$--$88$), Acceleration's count,
the slot balance, and the rank give on the order of $10^{11}$ combinations. One identity
collapses it. A target locked at count $c_i$ uses $(89-c_i)/3$ feeds, and since every
increment lands on Acceleration or on a target before it locks, $\sum_i c_i + a = 245$.
Hence the total feeds spent is
\begin{equation}\label{eq:budget}
   \sum_{i=1}^{4}\frac{89-c_i}{3}=\frac{356-\sum_i c_i}{3}=\frac{111+a}{3},
\end{equation}
\emph{independent of when the locks happen}. So the $50$-slot budget is met exactly when
$a\le 39$, the very condition the Acceleration cap already imposes. The feed count
therefore never needs tracking, and each target collapses to its residue $n_i\bmod 3$
(only ``on a window?'' matters). Dropping the secondary requirement that slots be
physically accrued by lock time (verified below to bind on just $0.17\%$ of runs), the
state shrinks to
\[
   (\text{rank},\ a,\ \text{multiset of unlocked residues}),
\]
fewer than $10^{5}$ values, an exact backward induction that runs in milliseconds.

\subsection{The optimal value}
Backward induction on this state yields the optimal value of the (slot-relaxed) online
problem,
\[
   \boxed{\,1.170\%\,}\ \ (\text{Acceleration}\le 250),\qquad 0.979\%\ \ (\text{Acceleration}=250).
\]
Two internal checks support it. Under a fixed feed-last policy the same transition model
reproduces the exact Grade-3-only value $0.10239\%$, and the optimum exceeds the best
fixed-deadline rule of Section~\ref{sec:policy}, as it must. Because the relaxation only
\emph{removes} a constraint, $1.170\%$ is an \emph{upper bound} on the true online
optimum. Replaying the optimal residue-policy with the slot-timing constraint reinstated
achieves $1.063\%$ (a Monte-Carlo estimate), a matching lower bound. Therefore
\[
   1.063\%\ \le\ (\text{true online optimum})\ \le\ 1.170\%.
\]

% =====================================================================
\section{The optimal policy is an Acceleration-aware threshold}\label{sec:policy}
% =====================================================================

The optimal action has a clean shape. \emph{Lock a target the moment it reaches a
window, provided $\text{rank}\ge T(a)$}, where the threshold $T$ rises with the
Acceleration count $a$ (Table~\ref{tab:thresh}). When Acceleration is low there is
budget to spare, so banking a window early pays off. When Acceleration is tight one
must wait until the final rank or two. A simple fit is $T(a)\approx 39+a/5$.

\begin{table}[h]
\centering
\begin{tabular}{@{}rrc@{}}
\toprule
Acceleration count $a$ & Acceleration value & lock from rank $T(a)$ \\
\midrule
0  & 55  & 39 \\
8  & 95  & 41 \\
16 & 135 & 42 \\
24 & 175 & 44 \\
30 & 205 & 45 \\
36 & 235 & 47 \\
39 & 250 & 49 \\
\bottomrule
\end{tabular}
\caption{Optimal lock threshold with four targets still unlocked. Feed a target to
$500$ on its divisible-by-$15$ window once $\text{rank}\ge T(a)$. With fewer targets
left, lock about one rank earlier.}
\label{tab:thresh}
\end{table}

This is why a \emph{fixed} deadline is suboptimal. The best fixed rule, ``never lock
before rank $D$, then lock on sight,'' peaks at $D=47$ and scores $0.885\%$
(Monte-Carlo). It is too patient in the low-Acceleration states, where the optimum
locks as early as rank $40$. The Acceleration-aware rule recovers that, raising
success by about $20\%$ ($0.885\%\to 1.063\%$), roughly one perfect bird in $90$
rather than one in $113$.

% =====================================================================
\section{The offline ceiling, bounded}\label{sec:ceiling}
% =====================================================================

The offline ceiling is the only figure in this paper not computed exactly, and it
is worth seeing why. It is an \emph{offline} quantity, the fraction of climbs for which
\emph{some} lock schedule wins \emph{given foreknowledge of every roll}. Backward
induction cannot reach it, because the optimal decision depends on the future and,
through exclusion, each decision reroutes that very future. So it is \emph{bracketed}
exactly, then pinned by simulation.

\paragraph{Exact upper bound.}
Remove the exclusion side effect, and let a locked target keep absorbing its own
increments. Then Acceleration keeps its natural count $A$, and every target passes
through a window (at counts $2,5,8,\dots$) for free, so success reduces to the budget
$A\le 39$ alone. Since this can only help, by Proposition~\ref{prop:naive}
\[
   P_{\text{clair}}\ \le\ \Prob(A\le 39)\ =\ 6.166840\%.
\]

\paragraph{Exact lower bound.}
Feed-last is a concrete feasible schedule, so its success rate is a rigorous floor.
That rate requires all four counts $\equiv 2\pmod 3$ and $A\le 39$, and equals
$0.10239\%$ (the value validated in Section~\ref{sec:mdp}). Combining,
\[
   0.10239\%\ \le\ P_{\text{clair}}\ \le\ 6.16684\%.
\]

\paragraph{Pinning the value.}
Inside this bracket the value is located by high-precision Monte Carlo. For each climb
an exhaustive lock-schedule search runs on its fixed future, exactly the best a
player who could see every roll would do. With $N=10^{6}$ runs (reproducible seed),
\[
   P_{\text{clair}} = 4.476\% \pm 0.04\%\qquad(\text{Wilson 95\% interval }[4.44\%,\,4.52\%]),
\]
comfortably inside the proven bracket. The two exact bounds together with this
confidence interval are the rigorous statement of the ceiling, and a closed form is not
available, for the reason above. The prophet gap is the value of foreknowledge, a
factor of roughly four between the online optimum ($\le 1.170\%$) and this ceiling, and
no online policy can recover it.

% =====================================================================
\section{The exact lineup: Acceleration at exactly 250}\label{sec:exact250}
% =====================================================================

The goal of Part~II so far asked only that Acceleration finish at \emph{or below}
$250$. A stricter reading insists on the literal lineup $500/500/500/500/250$, with
Acceleration at \emph{exactly} $250$ (count $a=39$). The numbers change sharply, and
one structural fact explains why.

\paragraph{At the exact lineup the feed rule is moot.}
If $a=39$ the four targets share $245-39=206$ increments, so by the budget identity
\eqref{eq:budget} their total deficit is $5(4\cdot 89-206)=750$ points. Closing $750$
points in $50$ slots forces an average of exactly $15$ points per slot, so
\emph{every} feed must be Grade-3, and each target's deficit must itself be a multiple
of $15$. Mixed feed grades buy nothing here. The exact lineup is intrinsically a
Grade-3-only target, and the ``$\le 250$'' slack that mixed feeds exploit is gone.

\paragraph{The probabilities.}
Pinning Acceleration to exactly $250$, the methods of
Sections~\ref{sec:mdp}--\ref{sec:ceiling} give the following, alongside the $0.979\%$
relaxed online optimum already recorded in Section~\ref{sec:mdp}.
\begin{itemize}[leftmargin=1.4em]
  \item \textbf{Feed-last} reaches $0.067\%$ (exact). It needs both $a=39$ and all
        four residues aligned, with no adaptivity to help.
  \item \textbf{Online optimum} lies between $0.847\%$ (achievable, real slot-timing,
        Monte-Carlo) and $0.979\%$ (relaxed, exact), the same lower and upper sandwich
        as the $\le 250$ case, only stricter.
  \item \textbf{Offline} reaches $4.36\%$ (Monte-Carlo, Wilson $95\%$ interval
        $[4.28\%,\,4.46\%]$).
\end{itemize}

\paragraph{Reading against the $2.30\%$ benchmark.}
It is tempting to compare the exact lineup to the $2.301513\%$ four-attribute figure
of Part~I, but they are different targets. The $2.30\%$ figure asks for four caps with
Acceleration anywhere $\le 250$ and \emph{any} feed grade, fed at the end. Pinning
Acceleration to exactly $250$ sits at the budget edge and discards that slack, so a
real-time player cannot approach $2.30\%$, and the online ceiling is about $1\%$. Only
an offline player, who can deliberately inflate Acceleration up to exactly $250$ by
locking targets early, exceeds it (at $\approx 4.36\%$), and that requires
foreknowledge, so it is not playable. In short, the literal $500/500/500/500/250$
lineup is a vanity constraint. Since Acceleration is the dump attribute, ending at
$230$ is as good as $250$, and the achievable goal worth pursuing remains the
$\le 250$ four-attribute target of Part~I.

% =====================================================================
\section{Results and discussion}\label{sec:results}
% =====================================================================

Table~\ref{tab:master} collects every probability computed in the paper, with the
method that produced it.

\begin{table}[h]
\centering
\begin{tabular}{@{}llcl@{}}
\toprule
Scenario & Play & Value & Method \\
\midrule
\multicolumn{4}{@{}l}{\itshape Fixed-end feeding (Part~I)}\\
Idealized bound      & frictionless $a\le 39$    & $6.166840\%$            & exact (binomial sum) \\
Four, fixed dump     & max $4$, Acceleration dumped & $2.301513\%$         & exact (gen.\ function) \\
Four, flexible dump  & max $4$, worst dumped     & $11.324144\%$           & exact (gen.\ function) \\
Three attributes     & max best $3$              & $100\%$                 & exact (closed form) \\
\midrule
\multicolumn{4}{@{}l}{\itshape Adaptive, Grade-3-only perfect, Acceleration $\le 250$ (Part~II)}\\
Feed-last            & lock only at rank $50$    & $0.10239\%$             & exact \\
Best fixed deadline  & lock from rank $47$       & $0.885\%$               & Monte-Carlo \\
Online optimum       & best state-based policy   & $1.063\%$--$1.170\%$    & MC (lower), exact (upper) \\
Offline ceiling      & foresee every roll        & $4.476\%$               & MC; bounds exact \\
\midrule
\multicolumn{4}{@{}l}{\itshape Adaptive, exact lineup $500/500/500/500/250$ (Part~II)}\\
Feed-last            & lock only at rank $50$    & $0.067\%$               & exact \\
Online optimum       & best state-based policy   & $0.847\%$--$0.979\%$    & MC (lower), exact (upper) \\
Offline              & foresee every roll        & $4.36\%$                & Monte-Carlo \\
\bottomrule
\end{tabular}
\caption{All results. The $\le 250$ offline ceiling is bracketed exactly within
$[0.10239\%,\,6.16684\%]$ and pinned by simulation at $4.476\%$ (Wilson $95\%$
interval $[4.44\%,4.52\%]$, $N=10^{6}$). The exact-lineup offline value is the
analogous Monte-Carlo estimate. All figures are exact except the adaptive-regime
Monte-Carlo values noted, namely the best fixed deadline, the online achievable lower
bounds, and the offline point values.}
\label{tab:master}
\end{table}

\paragraph{Reading the fixed-end numbers.}
The headline $2.301513\%$ means that a player who commits to dumping Acceleration and
feeds at the end perfects the other four attributes on about one final chocobo in
$1/0.0230\approx 43$. The shortfall from the idealized $6.17\%$ is entirely feed
chunkiness. A leftover of $5$ or $10$ points wastes a whole slot, so the real budget
is tighter than a points count suggests. Refusing to pre-commit, and dumping whichever
attribute rolls worst, raises the rate to $11.32\%$ (about one in nine), and merely
perfecting the best \emph{three} attributes is certain.

\paragraph{Reading the adaptive numbers.}
The strict Grade-3-only lineup $500/500/500/500/250$ is far rarer, but \emph{when}
the player feeds matters enormously. Feeding only at rank $50$ succeeds just
$0.10239\%$ of the time. An Acceleration-aware lock-timing policy lifts the online
optimum to between $1.063\%$ and $1.170\%$ (about one in $90$), roughly $20\%$ above
the best fixed deadline. An offline player who could foresee the rolls would reach
$4.476\%$, but the gap from the online optimum to that ceiling is the prophet gap,
unrecoverable without foreknowledge.

\paragraph{A rule of thumb for play.}
For the any-grade goals of Part~I, the advice is simply to bank all feeds and spend
them at rank $50$. One may also abandon a doomed bird early, since with $k$ rank-ups
left the four targets' combined deficit can fall by at most $40k$ points, so once it
exceeds $40k$ the bird cannot recover. For the Grade-3-only lineup of Part~II, watch
Acceleration. If its value passes $250$ the perfect finish is already impossible.
Otherwise feed a target to $500$ the moment its gap is a multiple of $15$ and the rank
has reached the threshold $T(a)$ of Table~\ref{tab:thresh}.

\paragraph{Exact versus simulation.}
Every fixed-end figure, the online optima's upper bounds, the feed-last values, and
the $\le 250$ offline ceiling's two bounds are exact. The remaining figures resist
closed form and are estimated by simulation. These are the best fixed deadline ($0.885\%$), the
online achievable lower bounds ($1.063\%$ for $\le 250$ and $0.847\%$ for the exact
lineup), and the offline point values ($4.476\%$ and $4.36\%$). Each is flagged in
Table~\ref{tab:master} and where it is derived.

% =====================================================================
\section{Conclusion}\label{sec:conclusion}
% =====================================================================

I modeled the climb of a final \emph{Final Fantasy XIV} race chocobo as a
multinomial allocation of $245$ random growth increments, refined by a $50$-slot
feeding budget delivering $5$, $10$, or $15$ points per slot, and studied the chance
of perfecting four attributes under two regimes of play.

In the \textbf{fixed-end} regime the player banks every feed and spends it at rank
$50$. I proved this schedule optimal for the any-grade goals, showed that the
familiar $6.17\%$ community figure answers only a frictionless idealization, and
computed the true probabilities once feeds are recognized as discrete chunks. These
are $2.301513\%$ for a pre-committed Acceleration dump, $11.324144\%$ when the dump is
chosen after the growth is seen, and certainty for the three-attribute goal. Each of
these is exact, by a binomial sum, a generating-function enumeration, and an
elementary inequality.

In the \textbf{adaptive} regime the player feeds during the climb, exploiting the rule
that a maxed attribute leaves the random pool. Casting the climb as a Markov decision
process and collapsing its state through the identity that total feeds depend only on
the final Acceleration count, exact backward induction brackets the best online
probability of the strict Grade-3-only perfect (four maxed, Acceleration $\le 250$)
between $1.06\%$ and $1.17\%$. That optimum is realized by an Acceleration-aware
lock-timing threshold, which locks a target on its divisible-by-$15$ window once the
rank reaches $T(a)$, and it beats any fixed deadline by about $20\%$. An offline
player who could foresee the rolls would reach $4.48\%$, exactly bracketed within
$[0.10\%,6.17\%]$, and the gap from the online optimum to that ceiling is the prophet
gap, the value of information no real player commands. Demanding the \emph{exact}
lineup $500/500/500/500/250$ is stricter still. At that budget edge the feed grade
becomes irrelevant (every feed must be Grade-3), the online ceiling falls to about
$1\%$ ($0.85$--$0.98\%$), and only an offline player, able to inflate Acceleration up
to exactly $250$, reaches $\approx 4.4\%$.

For a player, the practical reading is plain. Three capped attributes come for free,
four are a one-in-$43$ to one-in-nine lottery depending on how the dump is chosen, and
the Grade-3-only perfect (Acceleration $\le 250$) is about one bird in ninety even
under optimal timing. The literal $500/500/500/500/250$ lineup is rarer still and
not worth the fuss, since a lower Acceleration is just as good. All figures are exact
except the handful of adaptive-regime simulation estimates, flagged throughout.

\appendix
% =====================================================================
\section{Model constants}\label{app:constants}
% =====================================================================
\begin{center}
\begin{tabular}{@{}ll@{}}
\toprule
Quantity & Value \\
\midrule
Attribute cap (final chocobo) & $500$ points \\
Points per $1\%$ of cap & $5$ \\
Starting value per attribute (rank $1$) & $55$ points ($11\%$) \\
Rank-ups from rank $1$ to $50$ & $49$ \\
Random increments per rank-up & $5$ (each $+5$ points) \\
Total random increments $N$ & $245$ \\
Feeding slots over a lifetime & $50$ \\
Points per feeding (grades $1/2/3$) & $5\,/\,10\,/\,15$ \\
Increments to max one attribute by chance & $89$ \\
\bottomrule
\end{tabular}
\end{center}

% =====================================================================
\section{Full lock-threshold table}\label{app:thresh}
% =====================================================================
Table~\ref{tab:thresh} in the main text lists the optimal lock threshold $T(a)$ for
four still-unlocked targets at a coarse grid. Table~\ref{tab:thresh-full} gives the
threshold for every Acceleration count $a$ from $0$ to $39$ and for each number $u$ of
still-unlocked targets. A target on a divisible-by-$15$ window should be locked once
$\text{rank}\ge T$.

\begin{table}[H]
\centering
\footnotesize
\setlength{\tabcolsep}{5pt}
\renewcommand{\arraystretch}{0.92}
\begin{tabular}{@{}rr cccc@{}}
\toprule
$a$ & Accel.\ value & $u=4$ & $u=3$ & $u=2$ & $u=1$ \\
\midrule
0  & 55  & 39 & 39 & 40 & 42 \\
1  & 60  & 39 & 40 & 40 & 42 \\
2  & 65  & 40 & 40 & 41 & 42 \\
3  & 70  & 40 & 40 & 41 & 42 \\
4  & 75  & 40 & 40 & 41 & 42 \\
5  & 80  & 40 & 40 & 41 & 43 \\
6  & 85  & 40 & 41 & 41 & 43 \\
7  & 90  & 41 & 41 & 42 & 43 \\
8  & 95  & 41 & 41 & 42 & 43 \\
9  & 100 & 41 & 41 & 42 & 43 \\
10 & 105 & 41 & 41 & 42 & 44 \\
11 & 110 & 41 & 42 & 42 & 44 \\
12 & 115 & 42 & 42 & 43 & 44 \\
13 & 120 & 42 & 42 & 43 & 44 \\
14 & 125 & 42 & 42 & 43 & 44 \\
15 & 130 & 42 & 42 & 43 & 45 \\
16 & 135 & 42 & 43 & 43 & 45 \\
17 & 140 & 43 & 43 & 44 & 45 \\
18 & 145 & 43 & 44 & 44 & 45 \\
19 & 150 & 43 & 44 & 44 & 45 \\
20 & 155 & 43 & 44 & 44 & 46 \\
21 & 160 & 43 & 44 & 44 & 46 \\
22 & 165 & 44 & 44 & 45 & 46 \\
23 & 170 & 44 & 45 & 45 & 46 \\
24 & 175 & 44 & 45 & 45 & 46 \\
25 & 180 & 44 & 45 & 45 & 47 \\
26 & 185 & 44 & 45 & 45 & 47 \\
27 & 190 & 45 & 45 & 46 & 47 \\
28 & 195 & 45 & 46 & 46 & 47 \\
29 & 200 & 45 & 46 & 47 & 47 \\
30 & 205 & 45 & 46 & 47 & 48 \\
31 & 210 & 45 & 46 & 47 & 48 \\
32 & 215 & 46 & 46 & 47 & 48 \\
33 & 220 & 46 & 47 & 47 & 48 \\
34 & 225 & 47 & 47 & 48 & 48 \\
35 & 230 & 47 & 48 & 48 & 49 \\
36 & 235 & 47 & 48 & 48 & 49 \\
37 & 240 & 47 & 48 & 48 & 49 \\
38 & 245 & 47 & 48 & 48 & 49 \\
39 & 250 & 49 & 48 & 49 & 49 \\
\bottomrule
\end{tabular}
\caption{Optimal lock threshold $T(a,u)$ from the backward-induction policy, for
\emph{every} Acceleration count $a$ from $0$ to $39$. Lock a windowed target once
$\text{rank}\ge T$. The threshold rises with Acceleration $a$ (less budget to spare)
and falls slightly as fewer targets remain ($u$ smaller). The lone exception is the
degenerate boundary $a=39$, where Acceleration already sits at its $250$-point ceiling
and any further increment dooms the run, so the thresholds there are effectively moot.}
\label{tab:thresh-full}
\end{table}

% =====================================================================
\section{Reproducibility}\label{app:repro}
% =====================================================================
The exact results use only finite arithmetic. The binomial sum
\eqref{eq:binomcdf} is evaluated in exact integers. The chunked probabilities
(Section~\ref{sec:exact}) and the online optimum (Section~\ref{sec:mdp}) are
generating-function convolutions and backward inductions over fewer than $10^{5}$
states, carried with their internal checks (total probability $1$, and reproduction
of the idealized $6.166840\%$ and the feed-last $0.10239\%$). The three Monte-Carlo
figures use fixed seeds and the sample sizes quoted where they appear, most notably
$N=10^{6}$ for the offline ceiling, reported with a Wilson $95\%$ interval. No
result depends on unstated randomness.

% =====================================================================
\begin{thebibliography}{9}
% =====================================================================
\bibitem{cgw} \emph{Chocobo Racing}, Final Fantasy XIV Online Wiki (Console Games
Wiki). \url{https://ffxiv.consolegameswiki.com/wiki/Chocobo_Racing}

\bibitem{wp-rankup} \emph{Rank-Ups: Attribute Gains}, FFXIV Chocobo Racing.
\url{https://ffxivchocoboracing.wordpress.com/2015/06/01/intermediate-rank-ups-attribute-gains/}

\bibitem{wp-final} \emph{Advanced: Raising That Final Chocobo}, FFXIV Chocobo Racing.
\url{https://ffxivchocoboracing.wordpress.com/2015/06/02/advanced-raising-that-final-chocobo/}

\bibitem{wp-train} \emph{Training}, FFXIV Chocobo Racing.
\url{https://ffxivchocoboracing.wordpress.com/2015/05/26/intermediate-training/}

\bibitem{icy} \emph{Chocobo Racing and Breeding Guide}, Icy Veins.
\url{https://www.icy-veins.com/ffxiv/chocobo-racing-and-breeding-guide}

\bibitem{mdp} M.~L.~Puterman, \emph{Markov Decision Processes: Discrete Stochastic
Dynamic Programming}, Wiley, 1994.

\bibitem{prophet} T.~P.~Hill and R.~P.~Kertz, \emph{A survey of prophet inequalities
in optimal stopping theory}, Contemp.\ Math.\ \textbf{125} (1992), 191--207.
\end{thebibliography}

\end{document}
